\scanpage{93}

Bien entendu, on récupère par descente la filtration de $F$ connaissant $F'=F^P$, $i:G_m\longrightarrow G'=\underset{\underline{\mathrm{Aut}}_S(F')}{\underset{\shortparallel}{G^P}}$, et le sous-$H'$-torseur \inserted{à gauche} $Q$ de $P$.

\uuline{En résumé} (si $k$ de car. $0$, $\mathcal M$ $1/2$ simple), la donnée d'un \uline{foncteur} \inserted{\uline{fibre}} \uline{filtré} \uline{admissible} sur $\mathcal M$, à ``valeurs dans $S$'', équivaut à la donnée

(i) D'un foncteur fibre \uline{gradué} \uline{admissible} \inserted{$F'$} sur $\mathcal M$, à valeurs dans $S$, i.e. d'un foncteur fibre $F'$ sur $\mathcal M$, et d'un hom $i_1:G_{m_S}\longrightarrow G'=\underline{\mathrm{Aut}}_S F'$.

(ii) D'un $G'^{>}\deleted{\unsure{H'|G'^{\geq}}}$-torseur \deleted{$P$} $Q'$ à droite, où $\deleted{H'}G'^{>}$ \deleted{\unsure{est}} est le \deleted{\unsure{ss}} groupe uni-parabolique de $G'$ défini par $i$.

A ces données, on associe le foncteur fibre $F=F'^{Q'}$, avec la filtration déduite par descente de la filtration \deleted{canonique} de $F'$ \deleted{\unsure{gradué|qui lui}} associée \unsure{canon$^{\mathrm{t}}$} à sa graduation. D'ailleurs $Q$ est trivial sss la filtration du foncteur $F$ ``splitte'' $\ldots$

Si $\mathcal M$ est muni de\deleted{s} \inserted{la} donnée \unsure{n°} $1$ des structures \unsure{motiviques|motoriques|intrinsèques}, \deleted{\unsure{on}} alors on a $j:G_{m_S}\longrightarrow G'$ hom central, et à $i_1$ il convient de joindre $i_2$ défini par

\[
i_1(\lambda)i_2(\lambda)=j(\lambda)
\]

Alors $i_1,i_2$ sont dits \uline{complémentaires}. Les paraboliques associés sont opposés \deleted{(\unsure{$G^{>}$})}, car leur intersection est un Lévi commun.
% archive sheet number: 92
